\section{GlyphChain Legal Document Proof, Recovery and Evidence Packaging}
\label{sec:glyphchain-legal-document-proof}

\subsection{Purpose and Product Boundary}

Tessaris now provides a privacy-preserving evidence layer for documents completed
through the Legal workspace.  The feature binds an immutable document version,
the signing manifest and the envelope audit history to a portable proof receipt.
Where GlyphChain submission is enabled, the receipt also records the transaction
and its Merkle inclusion evidence.  A completed envelope can be exported as a
self-contained private evidence package for retention by the parties.

The design deliberately separates proof from publication.  Document text,
signer names, signer email addresses and the detailed audit history remain in the
originating Tessaris workspace and in the private package selected by the user.
GlyphChain receives only non-reversible commitments and protocol metadata.  The
feature does not depend on Photon, does not transfer currency and does not expose
private document contents on-chain.

\begin{quote}
\textbf{Current assurance statement:} Tessaris provides tamper-evident electronic
signing, a device-signed document commitment, recoverable proof submission and
portable verification evidence.  This is not a claim that Tessaris is a qualified
trust service provider, or that a signature is a qualified electronic signature
or qualified electronic seal.  Verification of integrity does not determine
whether a document is lawful, fair or enforceable.
\end{quote}

\subsection{Implemented Features and Functionality}

The completed local foundation combines a controlled Legal workspace, private
electronic-signing workflow, portable evidence packaging and a GlyphChain
integrity and finality rail. The principal user-visible capabilities are:

\begin{description}
  \item[Legal matters and evidence relationships.] Encrypted matter records can
        link protected documents, communications, signature envelopes and proof
        receipts without duplicating the underlying private content. Renewal,
        expiry, notice and review obligations support owners, due dates,
        pre-due reminders, completion and waiver states.
  \item[Secure document intake.] Tessaris accepts bounded PDF, DOCX, text and
        supported image files, hashes the exact original bytes, stores the
        original and derived text under workspace-bound AES--256--GCM
        encryption, and performs bounded local OCR where required. Unsafe,
        executable, mismatched or resource-excessive inputs are rejected.
  \item[Contract review and comparison.] Exact document versions can be
        reviewed for explainable clause and risk signals, compared using a
        bounded unified diff, and used to prepare non-destructive redline or
        negotiation drafts. Automated findings remain triage guidance requiring
        human review and are not represented as legal advice.
  \item[Legal Pilot.] The local Legal harness accepts protected document
        identifiers rather than raw desktop paths, preserves workspace
        isolation, produces evidence references and capability receipts, and
        cannot autonomously send, sign, approve or modify an external document.
  \item[Controlled template library.] Customer, supplier, contractor,
        employment, terms, privacy and data-processing workflow classes are
        supported through encrypted, immutable and versioned templates.
        Approval is tied to the exact content digest, named reviewer assertion,
        qualification, jurisdiction and review scope. Instantiation produces a
        final-review-required encrypted document; no generic wording is bundled
        or described as legally approved.
  \item[Electronic-signing workflow.] Draft, review, exact-version approval,
        ordered signing, optional separate-channel possession codes, completion,
        reminders and signed-copy evidence are implemented. Plaintext signer
        links and one-time codes are disclosed once; only their hashes are
        retained.
  \item[Proof cabinet and portable verification.] Completed envelopes expose
        the document digest, sealed manifest digest, transaction reference,
        Merkle inclusion evidence, quorum-certificate fields and append-only
        event timeline. Users can download private evidence archives and
        \texttt{.tessaris-proof} receipts and verify them locally without
        uploading the document.
  \item[Controlled signed-copy delivery.] An encrypted, restart-safe and
        idempotent outbox prepares exact evidence-package attachments, applies
        bounded retry, records provider outcomes, and requires reauthorisation
        following a terminal bounce or rejection. Live provider validation
        remains a deployment gate.
  \item[Records governance.] Retention deadlines and legal holds block
        destruction. Deletion requires an export, exact confirmation and an
        encrypted content-free destruction receipt.
  \item[Encrypted device recovery.] An operator can create one bounded,
        password-encrypted recovery package containing all four Legal stores,
        signing and delivery state, the proof outbox, GlyphChain runtime state,
        a consistent SQLite ledger snapshot, public trust material and
        encrypted Legal and network-identity key-recovery bundles. A clean-root
        restore authenticates every encrypted record and audit chain before it
        is reported as successful.
  \item[Explicit network participation.] Desktop installations can be assigned
        light, relay or validator roles. Relay and validator activation requires
        signed-desktop administration, exact confirmations and role-specific
        acknowledgements; validator activation also requires governance and a
        matching trusted checkpoint. Role changes require restart before taking
        effect.
  \item[Network operations.] Signed peer discovery, bounded peer memory,
        transport-key rotation and recovery, restart and catch-up recovery,
        outbound-only NAT relay transport, malformed-input rejection, ingress
        backpressure and privacy-safe local network monitoring are implemented.
        Monitoring raises local alerts for
        finality stalls, conflicting finalized history, extended rounds,
        elevated-node isolation, peer pressure, queue drops and certificate or
        trust-configuration faults.
  \item[Proof-preserving storage maintenance.] A light-node operator can export
        and independently verify a complete finality-anchored ledger archive,
        compact only a separate database copy, retain a canonical restart-state
        base and finalized tip, detect concurrent writes, and atomically install
        the compacted ledger. The operation is explicit, reasoned, restart-gated
        and never triggered automatically by storage pressure. Portable Legal
        receipts remain verifiable after historical rows leave the hot ledger.
\end{description}

Professional escalation is intentionally a controlled draft. Tessaris may
assemble the matter context, evidence references, urgency and requested
specialism, but it does not appoint or contact a lawyer without a separate
authorised delivery action.

\subsection{Technical Architecture and Trust Separation}

The implementation separates confidential business records, proof commitments,
network transport identity and validator-finality authority. These domains are
not interchangeable:

\begin{enumerate}
  \item The \emph{Legal record plane} stores matters, documents, templates,
        envelopes, obligations and delivery state under workspace-isolated
        authenticated encryption.
  \item The \emph{evidence plane} produces an exact signed-copy manifest,
        hash-chained audit events, a portable receipt and a private evidence
        archive containing the material needed by the parties.
  \item The \emph{GlyphChain ledger plane} receives only non-reversible
        commitments and protocol metadata. It does not receive document text,
        signer names, email addresses or the private evidence archive.
  \item The \emph{consensus-finality plane} binds the exact canonical ledger
        block, transaction-envelope Merkle root and validator-set checkpoint to
        a portable quorum certificate.
  \item The \emph{network plane} uses a separate versioned Ed25519 transport
        identity for peer HELLO possession proofs and signed message exchange.
        Validator-finality keys cannot be silently substituted with the
        transport key.
\end{enumerate}

The main persistent components are an encrypted Legal record store, durable
proof-submission outbox, append-only submission journal, encrypted delivery
outbox, retained finality receipts, signed discovery-record store, versioned
network-identity store, bounded peer registry and bounded aggregate network
health history. Atomic file replacement is used for durable metadata and
owner-restrictive permissions are applied to sensitive local files. Production
key custody uses the operating-system Keychain and fails closed if that custody
service is unavailable; deterministic development mode uses explicitly scoped
owner-only files.

\subsubsection{Cryptographic binding}

For document bytes $D$, ordered signer commitments $S$, the terminal audit
event $A$ and the sealed signed-copy manifest $M$, Tessaris computes

\[
 h_D=\operatorname{SHA256}(D),\qquad
 h_S=\operatorname{SHA256}(\operatorname{Canonical}(S)),
\]

\[
 h_A=\operatorname{SHA256}(A),\qquad
 h_M=\operatorname{SHA256}(\operatorname{Canonical}(M)).
\]

The \texttt{DOCUMENT\_PROOF\_V1} payload binds these values to a deterministic
proof identifier, origin-device key identifier, completion time and verification
policy. The receipt's transaction envelope is independently canonicalised and
hashed. Its inclusion path must reconstruct the declared transaction Merkle root.
The version-two quorum certificate then binds that same root and the canonical
ledger commitment to validator signatures checked against an independently
pinned validator-set checkpoint. A proof is not promoted from
\texttt{LEDGER\_COMMITTED\_AWAITING\_QUORUM\_FINALITY} to
\texttt{FINALIZED} unless these exact bindings verify.

\subsubsection{Idempotency, recovery and auditability}

Proof identifiers are deterministic per completed envelope. Before allocating
new chain state, Tessaris searches for an existing matching proof transaction.
Submission state is persisted before network activity and recovered after a
restart. Bounded exponential retry retains the attempt count, next-attempt time,
last error and update time; manual retry remains workspace scoped.

Legal storage and proof-signing keys are versioned and included in encrypted
recovery bundles. Network transport identities use an independent recovery
format bound to the exact chain and node. Rotation produces a continuity
certificate signed by both the retiring and replacement Ed25519 keys and retains
old public verification material. Recovery imports reject wrong passwords,
key/public-key mismatches, chain or node mismatches and invalid rotation chains.
The combined full-state recovery format described below now covers the local
Legal and GlyphChain device state. Application-account recovery,
deployment-managed validator-finality credentials and representative-device
production exercises remain separate gates.

\subsubsection{Encrypted full-state backup and clean-root restoration}

The \texttt{tessaris.legal-glyphchain-full-state.v1} package captures the four
Legal roots, including envelopes, matters, documents, templates, delivery
outboxes, proof receipts, proof retry state, destruction records and audit
history. It also captures GlyphChain consensus and participation state, a
SQLite online-backup snapshot of the ledger, the configured public validator
checkpoint and independently stored pin metadata. A canonical manifest records
the size and SHA--256 digest of every member, the source workspace, chain and
node, creation actor and time. The encrypted outer header intentionally omits
workspace, node and chain identifiers. The pin inside a recovery package is a
recovery copy only: it does not replace publication of the pin through an
independent channel and is never sufficient by itself to establish quorum
trust.

Raw Keychain items and development \texttt{\_secure\_keys} directories are not
copied into the archive. Legal stores use their existing password-wrapped,
versioned recovery bundles. The managed P2P identity uses its separate
chain-and-node-bound recovery bundle, including the verified rotation history.
A legacy proof-signing PEM is explicitly excluded and must first be migrated by
the Legal key store. Validator-finality private keys remain deployment-managed
secrets and are not scraped from process environment into a device backup.

The outer file uses scrypt with a unique 16-byte salt
(\(N=32768,r=8,p=1\)) and chunked AES--256--GCM. Each four-megabyte chunk has a
unique nonce derived from a random nonce prefix and monotonic counter; its
associated data binds the authenticated outer header, chunk index and exact
plaintext length. The header binds the total archive length, chunk count and
archive SHA--256 digest, detecting truncation, reordering, extension and partial
replacement. Export and restore require separate exact confirmation phrases and
a password of at least sixteen characters.

Before restoration, Tessaris authenticates every encrypted chunk, verifies the
archive digest, rejects absolute paths, parent traversal, links, special files,
duplicates, undeclared members, oversized files and excessive file counts, and
recomputes the canonical member manifest. All key-recovery bundles are decrypted
and validated in an isolated preflight root before destination state changes.
Restoration refuses a non-empty target unless the operator explicitly authorises
replacement; authorised replacement moves the previous state to a named,
recoverable safety path. Export start, export completion and restore completion
are appended to a hash-chained, fsync'd device-recovery journal. Installed
directories use owner-only permissions.

The clean-root regression restores and proves decryptability for every encrypted
record, verifies signing, matter and delivery event chains, recovers the proof
retry journal, checks SQLite integrity, preserves trust files and verifies the
network-identity rotation chain and public-key continuity. The operator command
\texttt{backend/scripts/legal\_full\_state\_recovery.py} exposes interactive
\texttt{export}, \texttt{inspect} and \texttt{restore} operations without
placing the recovery password on the command line. A restart remains mandatory
before network participation resumes.

The repeatable synthetic disaster-recovery exercise creates an encrypted
whole-device package, rejects a wrong password and a modified package, refuses
to overwrite an occupied target, restores all protected state into a clean
root, revalidates audit chains, the ledger, checkpoint and node identity, and
then runs the compromised/lost-validator governance exercise.  Its content-free
report records the exercise identifier, participants, topology, exact
implementation-file manifest, inputs, evidence digests, observed timings,
checks, findings and reproduction commands.  This is local technical evidence;
it is not application-account recovery, representative-device testing or
production approval.

\subsubsection{Application-account recovery authority}

Application-account ownership is now separated from restoration of local
device files.  Initialisation creates a stable account identifier derived from
an Ed25519 account-root public key, an encrypted recovery kit that must be kept
outside the Tessaris runtime, and a root-signed public account descriptor.  The
kit protects the root private key with scrypt
\((N=32768,r=8,p=1)\) and AES--256--GCM; its password is never retained.  Root
keys use platform Keychain custody in production.

Clean-device recovery requires the kit, password, expected account identifier,
current public descriptor, new device identifier and exact confirmation.  It
verifies root-key continuity, advances a monotonic recovery epoch, signs a new
descriptor, commits to the previous descriptor digest and revokes every device
not listed in the replacement descriptor.  Device authorisation fails for a
stale epoch or an omitted device.  Wrong passwords, modified kits or
descriptors, an unexpected account, an on-runtime kit destination and a
non-clean target all fail closed.

Where an account descriptor is present, the installed desktop request signer
and backend verifier now use the domain-separated
\texttt{aion.flow.desktop-session.v2} format.  The signed message binds the
account identifier and recovery epoch alongside the method, target, body
digest, device, person, workspace, timestamp and nonce.  A pre-recovery epoch,
a device absent from the active descriptor, a mismatched workspace or account,
and account headers received by a backend without the corresponding descriptor
all fail closed.  The legacy version~1 session format remains only as a
migration path when no account descriptor exists.

The interactive operator command
\texttt{backend/scripts/application\_account\_recovery.py} provides
\texttt{initialize}, \texttt{recover} and \texttt{status} operations.  It reads
recovery passwords from the terminal rather than command-line arguments,
writes only the root-signed public descriptor to a publication destination and
leaves the encrypted recovery kit under owner control.

The kit cannot self-assert that its supplied descriptor is the newest account
state.  The separately runnable registry implementation now authenticates the
root signature, preserves append-only history and enforces an exact monotonic
epoch and digest chain.  Its bounded public service stores no recovery secrets,
requires TLS for non-loopback use and is consumed by a fail-closed HTTPS client
with persistent rollback detection.  A co-located development instance is not
independent authority.  Production therefore still requires deployment under
separate operational control, connection of the remaining non-AION application
session issuers, and live clean-device recovery and lost-device revocation
validation against that registry.  Validator-finality credentials remain in
deployment-managed custody and are never included in the account recovery kit.

\subsubsection{Peer discovery, resource bounds and monitoring}

Peer bootstrap records are signed by independently pinned Ed25519 discovery
authorities and bind the chain, authority, monotonic sequence, issue and expiry
times, peer roles, endpoints and expected transport public keys. Nodes reject
unknown authorities, bad signatures, wrong-chain or expired records, excessive
lifetimes, replay, same-sequence equivocation, unsafe production endpoints and
oversized records. Signed-discovery enforcement disables legacy unsigned
connect and refresh routes. Installation of an expected key does not complete
trust: the live peer must still prove key possession through signed HELLO.

The peer registry defaults to a maximum of 2,048 entries, validates field
lengths and preferentially evicts the oldest non-validator record. Consensus
ingress uses bounded per-lane queues, duplicate suppression, response caching
for request-like lanes, per-peer message and byte rate limits, and explicit drop
reporting.

Persistent ledger growth is governed separately from peer memory. The effective
ledger limit is the lower of an explicit deployment limit and the storage limit
approved in the local network-participation record, with a 512 MiB default.
SQLite receives a hard main-database page ceiling, bounded WAL checkpoint policy
and journal-size target. Before a new block, transaction, checkpoint or metadata
write, Tessaris measures the database, WAL and shared-memory companions and
reserves configurable headroom. A write that would cross the writable limit
raises a dedicated capacity error before a new transaction is accepted into the
in-memory ledger. Idempotent retrieval of an existing transaction remains
available while new writes are blocked.
The storage administration route reports only byte counts, utilization and
policy state; it does not expose filesystem paths. Near-capacity and at-capacity
conditions feed the local health alert evaluator. Capacity exhaustion never
causes automatic record deletion: the policy is explicitly
\texttt{fail\_closed\_no\_automatic\_pruning}. Finality-anchored snapshot
synchronization and explicit proof-preserving pruning are implemented. Capacity
pressure alone never initiates deletion.

\subsubsection{Secure outbound-only NAT relay transport}

A desktop behind a home, office or virtual-machine firewall can receive
GlyphChain peer traffic without exposing an inbound port. It completes an
outbound HTTPS HELLO with an explicitly approved relay or validator, polls a
bounded recipient queue and passes every retrieved original P2P envelope
through its ordinary local router. The receiving node therefore rechecks the
sender's trusted public-key binding and message-lane signature. The relay is a
transport intermediary only: it cannot vote, appoint a validator, replace a
sender signature or create quorum finality.

Relay PUT, PULL and ACK operations use a separate canonical
\texttt{glyphchain.nat-relay.v1} request signed by the caller's Ed25519
transport identity. Each request binds the chain, operation, node identifier,
issuance time, unique nonce and complete operation payload. The relay requires
a previously verified HELLO, applies per-peer message and byte rate limits and
rejects stale, replayed, wrong-chain or invalidly signed requests. A target must
already have a public key in the trusted peer registry.

The restart-safe relay queue applies a global message bound, a per-recipient
bound, an exact encoded-message limit, a short maximum lifetime and bounded
pull batches. Identical envelopes are idempotent. Queue metadata is written
atomically under owner-only permissions. Recipients acknowledge an item only
after its local route accepted it or rejected it as permanently invalid;
transient server and rate-limit failures remain eligible for retry.

Signed discovery can bind a per-peer \texttt{relay\_url}. Consensus proposal
and vote broadcast selects that route without confusing the relay with the
destination node. A one-way HELLO response carries the responder's signed
identity, allowing mutual possession verification without an inbound callback.
When signed discovery is enforced, an unknown self-signed HELLO and any attempt
to replace a pinned peer key are rejected. HTTPS is mandatory outside an
explicit development-only insecure-network mode. Relay queues contain protocol
envelopes only and never contain Legal document text, signer identities or
private evidence packages.

\subsubsection{Finality-anchored snapshot synchronization}

A portable \texttt{glyphchain.ledger-snapshot.v1} archive contains exactly four
members: a consistent SQLite backup, a canonical manifest, the version-two
quorum-finality certificate and the public validator checkpoint. The manifest
binds the chain and network, snapshot height, state root, transaction root,
canonical ledger commitment, consensus block identifier, database byte length
and SHA--256 digest, certificate digest, checkpoint digest, checkpoint
identifier and validator-set root. Legal documents, signer identities and
private evidence packages are not included.

Verification is completed before installation. Tessaris rejects unexpected ZIP
members, traversal paths, encryption, excessive expanded size, database
corruption, foreign-key faults, unexpected tables, views or triggers, invalid
transaction envelopes, non-contiguous transaction ordering, checkpoint/root
mismatches, rollback below the local checkpoint, wrong-chain state and any
certificate that does not bind the reconstructed transaction root and ledger
commitment. A checkpoint carried in the archive cannot appoint itself: its exact
identifier and validator-set root must match the receiver's independently pinned
checkpoint.

Installation is restricted to an authenticated signed-desktop administrator,
an exact confirmation and effective light-node participation. The verified
database is copied to an owner-only staging file, hashed again, atomically moved
into place and strictly replayed. The prior ledger is retained as one local
recovery copy; a failed replay restores it. A restart is required before network
participation resumes.

For distributed synchronization, only an explicitly enabled relay or validator
with an active Ed25519 transport identity serves
\texttt{SNAPSHOT\_REQ}. Requests require an already HELLO-verified peer and a
signature over the chain-bound request payload. The receiving light node selects
an already verified peer by node identifier, streams the archive under its local
storage bound without redirects, verifies it using its own trust anchor and only
then invokes the same atomic installer. Stale snapshots and transport/archive
identity disagreements are rejected.

\subsubsection{Proof-preserving ledger pruning}

Historical compaction is exposed only as an authenticated signed-desktop
maintenance operation and is never automatic. The device must first be in its
non-relaying, non-voting light role and the operator must enter the exact
confirmation phrase and record a reason. Tessaris then exports the exact live
database into a complete finality-anchored snapshot, retains that archive with
owner-only permissions and independently verifies it against the operator's
pinned validator checkpoint.

The compactor works on a separate database copy. It binds a canonical chain-
state snapshot to the finalized height and state root, retains the finalized tip
block and its complete transaction envelopes, and removes only older block and
transaction rows. The state base allows deterministic restart without replaying
deleted history. New blocks use the maximum retained height plus one, rather
than assuming that retained rows begin at height one.

Before atomic installation, Tessaris produces a second consistent backup of the
live database and compares its digest with the retained archive. Any concurrent
ledger write aborts the operation. The generic atomic installer rehashes the
compacted database, retains the previous database for recovery, rebuilds runtime
state and verifies the finalized checkpoint before reporting success. Network
use requires restart.

Portable \texttt{.tessaris-proof} receipts do not depend on continued hot-
database retention: they carry the exact transaction envelope, transaction
hash, Merkle inclusion path, transaction root and quorum-finality certificate.
They therefore remain independently verifiable after compaction. The retained
full snapshot preserves historical inspection and recovery. This is not a claim
that deleting a local hot-ledger row deletes or revokes a distributed proof.
The archive directory may be placed on separately managed storage with
\texttt{GLYPHCHAIN\_PRUNING\_ARCHIVE\_DIR}; archive retention and backup-media
governance remain operator responsibilities.

Aggregate health snapshots contain no peer address, node identifier,
transaction, document or signature. Their bounded local hash chain is useful for
operator diagnostics but is not digitally signed and is not presented as an
independently authenticated forensic record.

\subsection{Implemented Interfaces}

The protected Legal and GlyphChain administration surfaces include:

\begin{itemize}
  \item document upload, classification confirmation, exact-version review,
        comparison, redline and negotiation workflows;
  \item template creation, revision, exact-digest approval, supersession and
        controlled instantiation;
  \item matter relationships, obligations, reminders and professional-
        escalation drafts;
  \item signature-envelope approval, signer actions, completion, proof,
        verification, evidence export and controlled delivery;
  \item proof-outbox status and workspace-scoped manual retry;
  \item light, relay and validator participation status, configuration and
        disablement;
  \item network-identity status, explicit rotation and encrypted recovery export
        and import;
  \item signed-desktop network-health capture and bounded history retrieval; and
  \item \texttt{POST /api/aion/business/glyphchain-admin/storage/status}
        for a path-free view of the effective ledger limit, database/WAL/shared-
        memory usage, remaining capacity, utilization and fail-closed policy;
  \item signed-desktop snapshot export, verify, explicit import and verified-peer
        synchronization routes;
  \item signed-desktop, light-node-only, explicitly confirmed archive-and-prune
        maintenance with concurrent-write detection and restart enforcement;
  \item authenticated \texttt{/api/p2p/nat/put},
        \texttt{/api/p2p/nat/pull} and \texttt{/api/p2p/nat/ack} routes for
        bounded outbound-only peer transport, plus aggregate path-free relay
        status; and
  \item the signed peer \texttt{POST /api/p2p/snapshot\_req} transport for a
        newer finalized snapshot from a relay or validator.
\end{itemize}

All administrative mutations are authenticated using a signed desktop session
bound to the HTTP method, request target and exact body. Workspace and operator
identity are derived from that verified session rather than accepted from an
untrusted request field.

\subsection{Completed Signing and Proof Flow}

The implemented flow is approval-gated and version-specific:

\begin{enumerate}
  \item Legal creates a draft envelope containing the exact document text and an
        ordered signer list.
  \item An authorised user submits the envelope for approval.  Tessaris creates a
        review hash over the envelope identifier, document hash, signer order and
        expiry time.
  \item Approval applies only when the supplied review hash still matches the
        version under review.
  \item Secure signer links are generated.  Plaintext tokens are returned once;
        only token hashes are retained locally.
  \item Each signer confirms the expected name and email address and accepts the
        exact document hash.  Ordered signing is enforced.  Per signer, the
        owner may additionally require a one-time code delivered through a
        separate channel.  Tessaris stores only the code hash and rotates both
        code and link when a reminder is prepared.  This is a possession check,
        not a claim of certified identity verification.
  \item When every signer has signed, Tessaris seals the signed-copy manifest and
        appends the completion events to the hash-chained envelope audit history.
  \item Tessaris creates a portable \texttt{.tessaris-proof} receipt and places
        its commitment in the durable GlyphChain proof outbox.
  \item If chain submission is enabled, the \texttt{DOCUMENT\_PROOF\_V1}
        transaction is submitted.  A temporary failure leaves the proof available
        for scheduled or manually requested retry.
  \item The user may download a private evidence archive containing the document,
        proof, certificate, audit events and verification instructions.
\end{enumerate}

\begin{verbatim}
Exact document + ordered signers
              |
              v
   Approval-bound envelope
              |
              v
 Signed-copy manifest + audit chain
              |
      +-------+--------------------+
      |                            |
      v                            v
Private evidence ZIP       Hash commitments only
                                   |
                                   v
                         Durable proof outbox
                                   |
                                   v
                     GlyphChain ledger commitment
                                   |
                                   v
                       Portable .tessaris-proof
                                   |
                        exact-block consensus binding
                                   |
                                   v
                    independently anchored quorum finality
\end{verbatim}

\subsection{Commitment and Receipt Model}

Let $D$ be the normalized UTF-8 document content and let

\[
  h_D = \operatorname{SHA256}(D).
\]

For every signer $i$, Tessaris computes a canonical commitment over the signer
identifier, signing time, signature method and accepted document hash.  The
canonical digest of the ordered commitment collection forms the signatures
commitment $h_S$.  The final event hash of the envelope's append-only audit chain
forms the audit commitment $h_A$.  The sealed signed-copy manifest has commitment
$h_M$.  The GlyphChain payload therefore binds:

\begin{itemize}
  \item a deterministic proof identifier;
  \item an opaque document-identity commitment;
  \item $h_D$, the document content hash;
  \item $h_M$, the signed-copy manifest hash;
  \item $h_S$, the ordered signature commitment;
  \item $h_A$, the terminal audit-event hash;
  \item the origin-device key identifier;
  \item the completion time, signing policy version and verification version.
\end{itemize}

The portable receipt schema is
\texttt{tessaris.glyphchain-document-receipt.v2}.  Its attestation comprises the
receipt version, creation time, normalized GlyphChain payload and origin-device
public-key description.  This attestation is signed with Ed25519.  The receipt
also contains an explicit privacy declaration and, after successful submission,
the transaction identifier, transaction hash, block height, transaction index,
transaction Merkle root and inclusion path.  A ledger commit is reported as
\texttt{LEDGER\_COMMITTED\_AWAITING\_QUORUM\_FINALITY}; it is not reported as
quorum-final unless the matching certificate verifies against a separately
pinned validator checkpoint.

The verifier recomputes the committed transaction envelope hash, checks that its
\texttt{DOCUMENT\_PROOF\_V1} payload exactly matches the receipt payload, checks
the transaction identifier, and then proves that the transaction hash belongs
to the declared Merkle root.  When finality becomes available, the receipt also
carries the version-two quorum certificate whose signed domain includes the
same transaction root and canonical ledger-block commitment.  Receipt contents
are never permitted to appoint their own trust anchor: the validator-set root
must come from a separately supplied checkpoint or operator-pinned local state.

\subsection{Durable Submission and Restart Recovery}

Proof submission is no longer dependent on an in-memory request or on the
desktop remaining open.  Each proof is written to a persistent outbox item before
submission.  Outbox state transitions are also written to an append-only,
filesystem-synchronized event journal.  The principal states are
\texttt{PENDING}, \texttt{SUBMITTING}, \texttt{RETRY\_PENDING},
\texttt{LEDGER\_COMMITTED\_AWAITING\_QUORUM\_FINALITY} and
\texttt{FINALIZED}.  The last state is reserved for a trust-anchored certificate
over the exact committed block.

Failed submissions use bounded exponential backoff.  Retry state includes the
attempt count, next-attempt time, most recent error and last update time.  A
manual retry API is available for controlled recovery, and the retry operation
is restricted to the authenticated workspace.  Proof identifiers are
deterministic for the completed envelope.  Before consuming a new chain nonce,
the service searches for an existing finalized \texttt{DOCUMENT\_PROOF\_V1}
transaction with the same proof identifier.  This provides proof-level
idempotency across request retries and application restarts.

The protected management routes are:

\begin{itemize}
  \item \texttt{POST /envelopes/\{id\}/proof} --- retrieve the private receipt;
  \item \texttt{POST /envelopes/\{id\}/proof/download} --- download the
        \texttt{.tessaris-proof} file;
  \item \texttt{POST /envelopes/\{id\}/proof/submit} --- submit or recover the
        proof on GlyphChain;
  \item \texttt{POST /envelopes/\{id\}/proof/verify} --- run server-side proof
        verification;
  \item \texttt{POST /proof-outbox/status} --- obtain workspace-scoped queue
        status;
  \item \texttt{POST /proof-outbox/retry} --- retry bounded workspace-scoped
        submissions; and
  \item \texttt{POST /envelopes/\{id\}/evidence-package/download} --- export the
        completed private evidence archive.
\end{itemize}

\subsection{Private Evidence Package}

The Legal interface exposes \emph{Download evidence} only for completed
envelopes.  The resulting \texttt{.tessaris-evidence.zip} archive is generated
in memory and is not written to GlyphChain.  It contains:

\begin{description}
  \item[\texttt{document.txt}] the exact frozen document content;
  \item[\texttt{signed-document.html}] a portable human-readable completed copy;
  \item[\texttt{evidence-certificate.json}] document identity, version, SHA-256
        hash, signer records, completion manifest, export time, proof reference
        and assurance limitations;
  \item[\texttt{audit-events.json}] the complete hash-chained envelope event
        history;
  \item[\texttt{document-history.json}] the verified replacement and
        supersession chain;
  \item[\texttt{\{envelope\}.tessaris-proof}] the portable device-signed proof
        receipt and any available GlyphChain inclusion evidence; and
  \item[\texttt{README.txt}] plain-language local verification and legal-status
        guidance.
\end{description}

The archive contains private document and identity information and must be
handled as confidential business evidence.  It is suitable for controlled local
retention or approved delivery to a party, but it must not be published or sent
automatically.  External signed-copy delivery remains disabled until an approved
connector and its authority boundary are configured.  Recording a delivery
reference does not itself claim that Tessaris sent the file.

\subsection{Approved Signed-Copy Delivery}

Completed evidence packages may now be prepared for Resend or
Microsoft~365/Graph delivery.  Preparation and transmission are separate
approval-gated actions.  Tessaris creates one encrypted outbox item for each
required signer, binds it to the completed manifest, and stores the exact ZIP
attachment before any provider call.  Deterministic delivery identifiers make
repeated preparation idempotent.

The outbox records four distinct facts rather than collapsing them into a
single ``sent'' state: approval and queueing, submission to the provider,
provider-reported delivery outcome, and confirmation for every required party.
A provider request identifier is mandatory before an item can become
\texttt{SUBMITTED\_TO\_PROVIDER}; that state is not displayed as delivered.
\texttt{DELIVERED}, \texttt{BOUNCED} and \texttt{REJECTED} require a separate
provider reference.  Failed submissions use bounded exponential retry and
survive process restart.  Bounced or rejected items retain their earlier
provider history and require an owner-authorised reason before another attempt.

Resend outcome ingestion is implemented as a public, provider-authenticated
callback.  Verification is performed over the exact raw HTTP body using the
\texttt{svix-id}, \texttt{svix-timestamp} and \texttt{svix-signature} headers and
one of the configured signing secrets.  Stale requests, invalid signatures,
oversized payloads and conflicting reuse of an event identifier fail closed.
The provider email identifier resolves through an encrypted opaque index to one
workspace-scoped outbox item; it is not accepted as a workspace identifier from
the callback.  Callback event identifiers are idempotent, while reuse with a
different payload digest is rejected.  Delivered, delayed, bounced, failed and
suppressed events update the outbox distinctly.  A contradictory final outcome
does not overwrite the earlier result: both facts are retained and the item is
marked for human review.  A verified delivered event is also appended to the
owning envelope's signed-copy delivery history.  Provider payload recipient and
message content are not copied into the callback journal.

The public callback path is
\texttt{/api/aion/business/legal-signing/webhooks/resend}.  Production requires
the \texttt{RESEND\_WEBHOOK\_SIGNING\_SECRETS} or singular
\texttt{RESEND\_WEBHOOK\_SIGNING\_SECRET} Vault-provided value and registration
of the exact HTTPS URL in Resend.  Multiple secrets are accepted during a
controlled rotation window.

Readable API responses omit the message body and attachment bytes.  Recipient,
message and evidence-package content remain inside the AES-256-GCM record, and
each state transition is protected by an internal hash chain.  Adapter limits
bound the number and aggregate size of attachments.  Production acceptance
still requires live Vault credential wiring and a real provider callback
exercise; tests do not send real email.

Microsoft~365 delivery uses a separate traceable flow.  Tessaris first creates
the message as a draft with Graph's immutable-identifier preference, records the
returned Internet Message-ID, and then sends that exact draft.  The Graph
\texttt{202 Accepted} response advances the item only to
\texttt{SUBMITTED\_TO\_PROVIDER}; it cannot confirm recipient delivery.  An
authorised refresh calls the Exchange Online Message Trace API and requires an
exact Internet Message-ID and recipient match before applying an outcome.  The
Graph status \texttt{delivered} becomes \texttt{DELIVERED}; \texttt{failed}
becomes \texttt{REJECTED}; and \texttt{quarantined} or
\texttt{filteredAsSpam} becomes \texttt{DELIVERY\_DELAYED} with human review.
Pending, expansion and unknown states remain non-delivered.  Trace observations
are bounded, idempotent, hash-journalled and contain no copied message body or
attachment.  A confirmed outcome is synchronized to the envelope's signed-copy
delivery history.

The traceable mail path requires the approved Microsoft~365 Vault connection to
grant \texttt{Mail.ReadWrite} for draft creation and \texttt{Mail.Send} for the
controlled send.  Message Trace additionally requires
\texttt{ExchangeMessageTrace.Read.All}, tenant administrator consent and the
Microsoft-documented Exchange tracing service-principal provisioning.  The
trace token may be supplied separately as
\texttt{MICROSOFT\_GRAPH\_MESSAGE\_TRACE\_ACCESS\_TOKEN}.  Local tests use
provider doubles; tenant onboarding, permission consent, trace latency,
throttling and real delivery outcomes remain live acceptance work.

\subsection{Independent Quorum-Certificate Verification}

The portable verification library accepts legacy version-one certificates for
compatibility, while the live content-bound path emits
\texttt{glyphchain.quorum-finality-certificate.v2}.  Independently of the live
consensus engine, it recomputes the canonical validator-set root, rejects
duplicate and non-member voters, verifies each Ed25519 precommit signature,
recomputes the two-thirds voting-power threshold, and checks the certificate's
chain, height, block, ledger commitment and transaction-root binding.  A
certificate is only reported as independently
trusted when its validator-set root matches a root supplied from trusted genesis
or checkpoint state.  Live validators retain a separate certificate-domain
Ed25519 precommit signature, assemble the quorum certificate, persist it in the
QC, reject unanchored QCs whenever finality assurance is configured, and expose
retained finality receipts.  Validator checkpoint files require an independently
pinned checkpoint identifier or validator-set root.  Validators independently
recompute every proposed transaction envelope, its transaction identifier, its
Merkle root and the canonical ledger commitment before voting.  The content
commitment is also embedded in the canonical consensus block identifier, so
votes carrying different ledger content cannot aggregate under one value.

The local browser verifier performs the equivalent checks without network
access.  It accepts a validator checkpoint as a separate fourth file and labels
a structurally valid certificate without that file as
\emph{signatures valid---trust anchor not supplied}.  The root copied into the
proof receipt is never treated as independent trust.  The Legal outbox fetches
the retained certificate after consensus and only promotes the receipt when its
transaction Merkle root matches the root certified by the quorum.

\subsection{Standalone Local Verification}

The static Tessaris GlyphChain Proof Verifier runs entirely inside the user's
browser.  Its content-security policy disables network, frame and external
resource access.  The selected document and receipt are not uploaded.  The
verifier currently checks:

\begin{enumerate}
  \item the SHA-256 hash of the normalized document against the receipt;
  \item the Ed25519 signature over the origin-device attestation;
  \item optional private signer commitments against the evidence certificate;
  \item the finalized transaction hash and exact proof-payload binding; and
  \item the transaction's Merkle inclusion path when committed evidence is
        available;
  \item validator membership, signed voting power and, when a separately
        supplied checkpoint matches, independently anchored quorum finality.
\end{enumerate}

The verifier reports each check independently.  Ledger commitment is not
treated as quorum finality, and a browser without Ed25519 WebCrypto
support reports the signature check as unsupported rather than silently passing
it.  The verifier is deliberately usable without a Tessaris account or network
connection.

\subsection{Consensus and Peer-Recovery Hardening}

The same delivery stage corrected the consensus failures exposed by the full
GlyphChain regression suite:

\begin{itemize}
  \item the peer router now imports and executes both vote and quorum-certificate
        ingress validation, preventing malformed quorum certificates from
        producing server errors;
  \item missing round fields use the protocol model's canonical default of round
        zero, aligning validation with accepted model construction;
  \item restart-sensitive synchronization refreshes the signed peer
        \texttt{HELLO} identity binding before requesting state or blocks;
  \item peer snapshots expose whether the signed identity handshake has
        completed;
  \item best-effort vote gossip reports observable bounded drops rather than
        generating retry-amplifying rate-limit failures; and
  \item dropped or full-lane work is explicitly reported as not accepted.
\end{itemize}

These changes close the previous gap-catch-up regression and stabilize restart
recovery, live finality, malformed-input handling, validation ordering and
backpressure behaviour.  The earlier intermediate result of 30 passing and nine
failing consensus tests is superseded by the final verified result below.

\subsection{Explicit Desktop Node Participation}

The versioned \texttt{glyphchain.network-participation.v1} record defines three
roles.  A light node verifies proofs and makes bounded outbound synchronization
requests but neither accepts public peers nor relays, proposes or votes.  A
relay requires an exact operator confirmation, a named operator, a public
connectivity acknowledgement and an explicit storage limit.  A validator also
requires availability and key-custody acknowledgements, a governance admission
identifier and a matching trusted validator checkpoint.

Role changes are atomic and appended to a hash-chained local audit journal.
When explicit enforcement is enabled, a missing record resolves to the light
role, and a validator-shaped environment variable alone cannot enable voting.
Light and relay roles suppress the local consensus validator identity.
Validator startup fails if the approved checkpoint does not match the trusted
checkpoint.  The Vault administration surface uses a signed, body-bound desktop
session and preserves all role-specific confirmations, acknowledgements and
checkpoint gates.  It writes the durable participation record but reports that
an application restart is required before capabilities change.  The production
installer default remains a deployment task; legacy developer clusters are
explicitly labelled rather than presented as the consumer security boundary.

\subsection{Matter Relationships, Obligations and Professional Escalation}

Legal matter records use the encrypted workspace store and now retain typed
references to protected documents, communications, signature envelopes and
proof receipts.  A relationship contains the protected record identifier, an
optional content hash and a human label; it does not duplicate the underlying
private document or message.  Duplicate relationships are rejected and each
change is appended to the matter's hash-chained audit history.

Renewal, expiry, notice, review and general-action obligations have explicit
UTC due times, optional pre-due reminder times, owners and source references.
Reminder times at or after the due time are rejected.  A workspace-scoped due
query reports triggered and overdue open obligations, while completed and
waived obligations are excluded.  The desktop Legal workspace provides the
same controls and exposes completion status without bypassing the encrypted
record service.

An escalation to a human legal professional creates only a structured draft.
It includes an urgency, requested specialism and a hash of the relevant matter
snapshot, but sets external delivery authority to false.  Tessaris does not
select, appoint or contact a lawyer automatically; a separately authorised
delivery workflow remains necessary.

The desktop signing workspace provides a proof-cabinet view for each envelope.
It displays the exact document digest, sealed manifest digest, GlyphChain
transaction, ledger commitment, transaction root, certificate schema and the
hash-chained event history.  Completed envelopes can export their private
evidence package from this view, and finalized proofs can invoke verification.
The view expressly distinguishes locally retained evidence from quorum finality
verified against an independently trusted checkpoint.

\subsection{Private Document Intake, OCR and Contract Review}

The Legal library accepts only an allowlist of bounded text, PDF, DOCX, PNG,
JPEG and TIFF inputs.  Active or executable formats, mismatched extensions,
empty files, oversized files, excessive PDF page counts, decompression-heavy
DOCX packages and excessive image dimensions are rejected.  The exact original
bytes are hashed with SHA--256 and stored with the derived text in a
workspace-bound AES--256--GCM record; plaintext originals are not retained by
the service.  List operations omit both original bytes and extracted text.

Text-layer PDFs, DOCX and UTF--8 text are extracted locally.  Supported images
and image-based PDFs use a bounded local Tesseract process when installed.
Page, embedded-image, pixel, execution-time and extracted-character limits
prevent unbounded document processing.  Missing text or unavailable OCR is
reported explicitly as \texttt{OCR\_REQUIRED} or
\texttt{OCR\_UNAVAILABLE}; Tessaris does not fabricate document content.

Document-category suggestions are explainable term matches and always carry
\texttt{REVIEW\_REQUIRED}, \texttt{human\_confirmation\_required=true} and
\texttt{legal\_determination=false}.  Confirmation records the responsible
person and time.  The document encryption-key domain is included in combined
Legal key recovery and rotation; older version-one recovery bundles remain
importable but explicitly report that they predate document-key custody.

The local contract-review engine binds every run to the exact document and
extracted-text digests.  It emits bounded clause-heading candidates,
source-excerpted rule signals and review guidance for matters such as automatic
renewal, unilateral termination, unlimited liability, indemnity, exclusivity,
restrictive covenants, personal-data duties, extended payment terms and an
unlocated governing-law term.  These are triage signals, not legal conclusions
or legal advice, and remain subject to authorised human or qualified
professional review.

Exact-version comparison records the two document digests and a bounded unified
line diff.  A redline proposal must quote text present in the exact source
document and is stored as a non-destructive draft; it does not alter the
original and carries no external-delivery authority.  Negotiation positions
record the issue, business position and counterparty position as locally
controlled open records.  The desktop review panel exposes classification,
local review, comparison, redline and negotiation controls under these same
boundaries.

The Legal Pilot chat route uses a versioned local harness rather than passing
raw desktop file paths into a generic conversation.  Its attachment contract
accepts only protected \texttt{document-*} identifiers, resolves them inside
the authenticated workspace, and rejects cross-workspace identifiers and raw
paths.  A review request invokes the exact-digest contract-review capability;
two attached versions can invoke exact comparison.  Each response includes
bounded evidence references and capability receipts and fixes
\texttt{external\_writes\_performed=false},
\texttt{external\_delivery\_authorized=false},
\texttt{human\_review\_required=true} and
\texttt{not\_legal\_advice=true}.  The bounded harness evaluation covers secure
attachments, review, comparison, workspace isolation, evidence-only overview
and external-authority denial.  It is not an evaluation of jurisdictional legal
correctness.

\subsection{Controlled Legal Template Library}

The Legal department includes an encrypted, workspace-isolated and versioned
template registry. It supports customer, supplier and contractor agreements,
employment documents, terms and conditions, privacy notices and data-processing
agreements as workflow classes. These classes are not bundled legal wording and
do not imply that Tessaris has approved a document.

Every template starts in
\texttt{DRAFT\_REQUIRES\_QUALIFIED\_REVIEW}. Approval requires a named reviewer,
their asserted qualification, the jurisdiction, the review scope, a positive
qualified-review acknowledgement and the SHA-256 digest of the exact content.
The system rejects a stale or mismatched digest. The approval record states that
Tessaris is recording the user's review assertion rather than certifying the
reviewer's status or providing legal advice.

An approved version is immutable. Editing creates a new version in the same
template family; approval of that version supersedes, but does not alter or
delete, the previous one. The audit history is itself hash chained. Only a
currently approved exact version may be instantiated. Placeholders use a
bounded \texttt{\{\{VARIABLE\_NAME\}\}} grammar; unknown and missing variables
are rejected. Instantiation creates an encrypted Legal document bound to the
template version, template digest and approval identifier. The resulting
document remains a
\texttt{CONTROLLED\_DRAFT\_REQUIRES\_FINAL\_HUMAN\_REVIEW}, carries no external
delivery authority and may optionally be linked to a Legal matter.

Template encryption keys participate in combined Legal key recovery and
rotation. The template library intentionally begins empty. Jurisdiction-specific
wording must be supplied and approved by a suitably qualified reviewer before
the interface permits controlled use.

\subsection{Signed Peer Discovery and Bounded Peer State}

Production discovery records use
\texttt{glyphchain.peer-discovery.v1}. Each record is bound to the chain,
authority identifier, monotonically increasing sequence, issue and expiry time,
and a bounded list of peer node identifiers, roles, endpoints and Ed25519
transport public keys. The canonical record is signed by an independently
pinned discovery authority. Nodes reject unknown authorities, invalid
signatures, wrong-chain records, expired or excessively long-lived records,
sequence replay, same-sequence equivocation, duplicate node identifiers,
oversized records and unsafe endpoint forms. Production policy requires HTTPS
and rejects loopback, link-local and private IP endpoints; a separate explicit
development setting permits local test nodes.

Only the newest accepted record per authority is retained using atomic local
replacement. The signed-discovery mode disables the legacy unsigned connect
and peer-list refresh routes. Remote refresh is restricted to an already known
peer, and every returned record is reverified locally before installation.
Discovery installs the expected transport key but does not mark the peer's live
HELLO challenge as complete; possession is still proven through the signed
HELLO exchange.

The runtime peer registry has a configurable hard count bound, rejects
unbounded identity fields and removes stale value-to-node aliases when a peer is
evicted. It prefers evicting the oldest non-validator peer before a validator.
Together with the fail-closed SQLite policy described above, this bounds peer
memory and persistent ledger growth. Finality-anchored snapshot transfer is
implemented as described above, together with explicit proof-preserving
pruning. Neither operation is triggered automatically by capacity pressure.

\subsection{Network Identity Rotation and Recovery}

The managed P2P transport identity uses versioned Ed25519 keys. Production key
material is held in the operating-system Keychain and fails closed when that
custody service is unavailable; deterministic development uses owner-only local
files. The router continues to accept an explicitly supplied legacy deployment
key, but otherwise reads the active managed identity for signed HELLO exchange.

Rotation requires a signed desktop administrator session, an exact confirmation
phrase and a recorded reason. A new private key is generated locally. The
rotation certificate binds the chain, node, sequence, old and new versions,
both public keys, actor, reason, timestamp and previous certificate hash. It is
signed by both the retiring and replacement private keys, creating verifiable
continuity rather than silently replacing a peer identity. Retired keys remain
available for historical verification and activation of a replacement key
requires an application restart.

All retained transport-key versions and the public rotation history can be
exported in an AES-256-GCM recovery envelope. Its key is derived with
scrypt using a unique salt; the envelope is bound to the exact chain and node.
Import requires explicit device-key-restore acknowledgement, rejects wrong
passwords, key/public-key mismatches, identity mismatches and invalid rotation
chains, and never returns plaintext private keys through the API. This mechanism
covers the network identity itself. The combined encrypted full-state package
adds the Legal records, queues, ledger, public trust material and audit history
around it. Application-account recovery and externally managed validator-
finality secret recovery remain open production gates.

\subsection{Privacy-Safe Network Health Monitoring}

GlyphChain exposes a signed-desktop administration route that captures a local
operator health snapshot. The snapshot contains aggregate consensus progress,
finality age and latency, peer counts by state, bounded-ingress queue totals and
operator alerts. It deliberately excludes peer URLs and IP addresses, node and
validator identifiers, transaction content, document hashes, signatures and
Legal workspace data. The chain name is represented only by a SHA-256 digest.
No snapshot is transmitted externally by the monitoring component.

Alert evaluation covers conflicting finalized block identifiers at the same
height, excess finality lag, extended consensus rounds, stale finality,
elevated relay or validator nodes with no verified peers, a peer registry near
its configured bound, new full-queue ingress drops, missing portable finality
certificates and finality trust-configuration errors. Alert thresholds are
locally configurable. This monitoring identifies conditions requiring operator
attention; it does not claim that a single local observation proves malicious
behaviour or establishes a production availability service level.

Snapshots are retained in a configurable bounded JSON-lines history. Every
record binds the previous record hash and its own canonical SHA-256 digest, so
accidental alteration and alteration without recomputing the retained chain are
detectable. The history is not digitally signed and is not represented as an
independently authenticated forensic record. When retention reaches its bound, the oldest
records are removed while the retained segment preserves its predecessor
commitment. Access to capture and history routes requires the enrolled signed
desktop administrator session.

\subsection{Fail-Closed Network-Partition Recovery}

GlyphChain recovery does not equate a peer's advertised height with trusted
finality.  Each synchronization cycle first validates every available peer
finality view against the independently configured validator checkpoint.  Peer
count is never substituted for certificate quorum power.  A higher valid view
is applied only after the node possesses and verifies the exact ledger blocks,
transaction root and ledger commitment covered by the certificate.

For a bounded gap, the node enters \texttt{BLOCK\_CATCHUP}, fetches the missing
range and verifies it before advancing persistent consensus state.  A gap above
the configured block-catch-up bound enters
\texttt{SNAPSHOT\_REQUIRED}.  In that case an administrator must use the
existing light-node snapshot workflow; the receiver verifies the archive,
public checkpoint and retained quorum certificate, atomically installs the
database, verifies the installed tip against that certificate, and only then
updates finalized consensus state.  This separates data installation from the
authority to finalize it.

Two cryptographically valid finalized views naming different blocks at the
same height, or a peer view conflicting with local finalized history, cause a
durable \texttt{SAFETY\_HALT}.  While halted, the engine may collect recovery
evidence but does not propose or vote.  The halt survives restart and ordinary
synchronization.  Its chain-bound and node-bound state is atomically stored
with owner-only permissions; a bounded hash-chained journal records detection,
phase transitions and clearance.

Clearance is available only to an enrolled signed desktop administrator and
requires the current halt identifier, a named operator, a reason and the exact
phrase \texttt{CLEAR GLYPHCHAIN PARTITION SAFETY HALT}.  Clearance records an
operator decision after the trusted history has been repaired; it cannot bypass
certificate, checkpoint, block or commitment validation.  Privacy-safe network
health exposes the recovery phase and aggregate progress but omits peer
identifiers, halt identifiers, transactions and document material.

\subsection{Protocol Upgrade Coordination}

Future consensus changes are coordinated through a durable activation plan
rather than inferred from an application version string.  The canonical plan
binds the current and target protocol versions, activation height, exact
release-manifest SHA-256 digest, trusted validator checkpoint and governance
approval.  The state is chain-bound and node-bound, atomically replaced with
owner-only permissions, and accompanied by a hash-chained operator journal.

The node advertises its active protocol, supported range and exact plan
identifier inside the Ed25519-signed HELLO identity domain.  Readiness is
therefore attributed only to an authenticated validator and measured by
validator voting power rather than by peer count.  Before activation,
insufficient quorum readiness is an operator warning.  At the activation
height, the consensus engine suppresses proposals, incoming votes and outgoing
votes unless the local release supports the target and matching signed
advertisements reach quorum power.  Unreadable or corrupt coordination state
also produces a fail-closed \texttt{UPGRADE\_HALT}.

Signed desktop administration controls scheduling, pre-activation cancellation
and post-activation completion.  Each operation uses an exact confirmation
phrase; scheduling additionally requires the release digest, checkpoint and
governance approval.  An activated plan cannot be cancelled locally.  After
the activation height finalizes on supported software, completion commits the
new current version and permits a later plan without discarding the audit
history.  These controls provide the upgrade rail; an actual future protocol
implementation and its representative rollout remain independently tested
release work.

\subsection{Operational Artefacts and End-to-End User Journeys}

The implementation is accompanied by explicit service-commitment, assurance,
privacy, network-participation and relay/validator operating records.  These
documents distinguish implemented safeguards from deployment promises.  In
particular, no numeric uptime, proof-submission, delivery or finality service
level is claimed before representative deployment and sustained-load evidence
exist.  The operator runbook covers role selection, independent trust-anchor
acquisition, key custody, explicit participation, pre-start checks, normal
operation, incident handling and retirement.  The runbook never treats message
submission as proof of delivery, ledger inclusion as quorum finality, or a hash
commitment as a recoverable copy of a private document.

The completed local product supports four principal end-to-end journeys:

\begin{description}
  \item[Review and negotiation.] A user imports an exact file, receives bounded
        local text extraction or OCR, confirms its classification, runs
        explainable contract triage, compares exact versions and records redline
        or negotiation drafts.  The original remains immutable and no draft is
        sent externally without a separate authorised action.
  \item[Agreement, signature and evidence.] An authorised person approves the
        exact document version and ordered signer list, issues one-time signer
        links, optionally applies a separate-channel possession code, completes
        the signing sequence, seals the signed-copy manifest, submits
        commitment-only evidence to GlyphChain and exports a confidential
        evidence archive and portable receipt.
  \item[Reusable controlled documents.] A business creates its own encrypted
        template, records qualified review against the exact content digest and
        jurisdiction, instantiates a final-review-required document and later
        creates a superseding version without overwriting the approved history.
  \item[Private desktop network operation.] An administrator keeps the ordinary
        installation in light mode or explicitly enables relay or validator
        participation.  Signed discovery, HELLO possession proofs, bounded peer
        state, outbound-only relay transport, independently pinned finality
        trust, local health alerts, verified snapshots and explicit archival or
        pruning preserve the privacy and resource boundary of the desktop.
\end{description}

Private originals, extracted text, signer identities, matter content, templates,
audit detail and evidence packages remain in the encrypted workspace unless the
user deliberately exports or delivers them.  The distributed rail receives
only canonical commitments and protocol metadata.  Public checkpoint
publication, provider delivery callbacks, independently operated relays,
representative-device recovery exercises and third-party security review
therefore remain release activities rather than hidden dependencies of the
local workflow.

\subsection{Verification Record}

The completed implementation has the following verified state:

\begin{itemize}
  \item \textbf{158 of 158} tests pass in the complete GlyphChain suite,
        including four-node finality, restart and gap catch-up, malformed quorum
        certificates, ledger-content commitments, finality-anchored snapshot
        export, verification and peer synchronization, proof-preserving ledger
        pruning and pruned restart continuity, signed outbound-only NAT relay
        authentication, bounds, delivery and discovery routing, backpressure,
        signature tampering, fail-closed network-partition recovery, verified
        large-gap snapshot adoption, sticky conflict halts, signed protocol
        readiness, activation-boundary halts and adversarial soak;
  \item the retained Legal-related tests cover secure storage, key recovery,
        encrypted full-state backup and clean-root restoration and its repeatable
        disaster exercise, off-device application-account recovery authority and
        signed desktop-session epoch enforcement, the monotonic standalone
        descriptor registry and fail-closed HTTPS registry client, matters,
        documents, delivery callbacks, signing, portable verification, the
        specialist workspace, controlled template review/withdrawal, and Resend
        and Microsoft Graph delivery adapters;
  \item the authoritative combined test count, duration, failures and errors are
        recorded in the current signed baseline report rather than duplicated as
        static prose in this section;
  \item the signed baseline contains the complete ordered path-and-SHA--256 source
        inventory, its canonical aggregate digest and file count, including the
        core release evaluator, release tooling, tests and every
        \texttt{GLYPHCHAIN*.md} and \texttt{GLYPHCHAIN*.tex} technical document;
        it verifies against the separately stored Ed25519 public trust anchor;
  \item the local resource-exhaustion review records request, signer-token,
        workspace-storage, record-enumeration, proof-queue, consensus-recovery
        and retained-state bounds; its focused regression set passes
        \textbf{93 of 93} tests;
  \item the retained compromised-validator exercise proves that an accused
        validator and a joining replacement cannot contribute approval power,
        while unaffected current-validator quorum produces a portable verified
        transition proof that survives governance-store restart;
  \item the repeatable disaster-recovery exercise passes all 25 checks for
        authenticated export, tamper and overwrite rejection, clean-root Legal,
        ledger, checkpoint and node-identity restoration, and lost-validator
        governance while retaining the application-account and production-test
        claim boundaries;
  \item the scheduled template lifecycle fails closed at the review deadline,
        permits exact-hash qualified renewal, and preserves approval history
        when a Legal change withdraws a version from controlled use;
  \item the cross-system incident-response playbook now defines containment,
        evidence classes, chain of custody, scenario recovery and return-to-
        service gates without treating a local backup as account ownership;
  \item Python syntax compilation passes for the signing, administration,
        snapshot, ledger and pruning components;
  \item JavaScript syntax validation passes for the Legal workspace and Vault
        administration controls; and
  \item the focused test structure was corrected so chain finalization and
        evidence export are independently exercised and reported.
\end{itemize}

The earlier 26-test targeted result and the temporary 30-pass, nine-failure
full-suite result are retained as development history only.  They do not
describe the current implementation state.  Focused and full-suite totals must
not be added together as a count of unique tests because the selections overlap.

\subsection{Security Properties and Present Limitations}

The implementation presently provides the following controls:

\begin{itemize}
  \item immutable document-version hashing and exact-version approval;
  \item one-time plaintext signer-link disclosure with only token hashes stored;
  \item restart-safe replay protection for signed desktop Legal and GlyphChain
        administration: HMAC-SHA-256 binds the method, exact target, body hash,
        device, person, workspace, timestamp and nonce, while a private bounded
        SQLite journal retains only device--nonce digests and fails closed if
        uniqueness cannot be established;
  \item ordered signers, explicit identity-field matching and optional
        separate-channel one-time possession codes;
  \item hash-chained envelope events and a sealed signed-copy manifest;
  \item Ed25519 device attestation;
  \item separation of Legal proof, P2P transport and validator-finality keys,
        with validator startup rejecting reuse of its transport key for
        certificate-domain votes;
  \item commitment-only chain payloads with no signer identity or document text;
  \item AES-256-GCM encryption for Legal matters and signing envelopes, with
        workspace, record, purpose and key-version context authenticated as
        associated data;
  \item production fail-closed custody in the operating-system Keychain and
        owner-only development key files;
  \item versioned storage and proof-signing keys, password-encrypted recovery
        bundles and verification of receipts made with retired signing keys;
  \item atomic proof, outbox, encrypted-record and key-metadata replacement;
  \item restrictive local file permissions for proof, record and key material;
  \item idempotent chain lookup and durable retry;
  \item restart-safe NAT relay nonce rejection, signed peer-discovery sequence
        control, durable delivery-callback event identifiers and terminal-state
        signer credential invalidation;
  \item workspace-scoped proof status and retry operations;
  \item materialize-before-finality partition recovery, bounded block catch-up,
        verified snapshot adoption and a restart-persistent conflict safety
        halt that suppresses validator proposals and votes;
  \item release-digest and checkpoint-bound protocol activation plans, signed
        validator readiness, quorum-power assessment and an activation-boundary
        halt for unsupported or insufficiently coordinated upgrades;
  \item bounded SQLite database, WAL and shared-memory growth with reserved
        write headroom, hard page limits and fail-closed transaction admission;
  \item legal holds and retention deadlines that block record destruction; and
  \item mandatory pre-deletion export, exact deletion confirmation and an
        encrypted content-free destruction receipt.
\end{itemize}

Production mode requires operating-system Keychain access and fails closed if
that custody service is unavailable.  The development and test boundary uses
owner-only files so automated tests remain deterministic.  The recovery bundle
uses a password-derived key and authenticated encryption; it contains no raw
plaintext keys.  Hardware-backed or secure-enclave custody remains a possible
future strengthening rather than a current claim.

\subsection{Remaining Production Gates}

The following items remain outstanding and must not be implied by the current
interface or marketing language:

\begin{enumerate}
  \item publish the exported trusted validator checkpoint file and its expected
        identifier through independent production release channels;
  \item register and validate the implemented Resend callback using live
        Vault-held credentials, and configure and validate the implemented
        Microsoft~365 traceable draft/send and Exchange Message Trace path with
        the required tenant permissions; Graph acceptance must not be
        represented as recipient delivery;
  \item complete application-account recovery, compromise-response runbooks and
        exercise the implemented restore procedure on representative production
        devices;
  \item deploy independently operated public relays and complete validator
        governance plus representative partition and future-upgrade exercises;
        signed discovery, outbound-only NAT transport, bounded peer state, role
        definitions, explicit local opt-in, fail-closed partition recovery and
        protocol-upgrade coordination are implemented;
  \item complete independent cryptographic, application-security, privacy and
        adversarial review; and
  \item assess any future qualified-signature or trust-service proposition with
        specialist legal and conformity advice before making regulated claims.
\end{enumerate}

GlyphChain is therefore a live, privacy-preserving document-proof rail inside
Tessaris, while Photon remains a separate future credit or compute economy.  No
Photon balance, token purchase or public document storage is required to create,
verify, recover or export the present Legal evidence.
